\documentclass[12pt]{article}
\usepackage{amsmath,amssymb,amsfonts}
\usepackage[margin=1in]{geometry}

\begin{document}

\section*{Problem 5.18}

\textbf{Problem:} $H_n(x)$ are the Hermite polynomials. Prove for arbitrary $x$:
\begin{itemize}
\item[(a)] $H_{2n}(0) = (-1)^n \dfrac{(2n)!}{n!}$
\item[(b)] $H_{2n+1}(0) = 0$
\item[(c)] $H_n'(0) = 2nH_{n-1}(0)$
\end{itemize}

\bigskip
\textbf{Solution:}

The Hermite polynomials are defined by the generating function:
\[
g(x,t) = e^{2xt - t^2} = \sum_{n=0}^{\infty} H_n(x)\frac{t^n}{n!}. \tag{$*$}
\]

\subsection*{Part (b): $H_{2n+1}(0) = 0$}

Setting $x = 0$ in the generating function:
\[
e^{-t^2} = \sum_{n=0}^{\infty} H_n(0)\frac{t^n}{n!}.
\]

Since $e^{-t^2}$ is an even function of $t$, all odd powers of $t$ must vanish. Therefore $H_n(0) = 0$ for all odd $n$, i.e.:
\[
\boxed{H_{2n+1}(0) = 0.}
\]

\subsection*{Part (a): $H_{2n}(0) = (-1)^n(2n)!/n!$}

From $e^{-t^2} = \sum_{n=0}^{\infty} H_n(0)\frac{t^n}{n!}$, using only even terms:
\[
e^{-t^2} = \sum_{n=0}^{\infty} H_{2n}(0)\frac{t^{2n}}{(2n)!}.
\]

But also:
\[
e^{-t^2} = \sum_{n=0}^{\infty}\frac{(-1)^n t^{2n}}{n!}.
\]

Comparing coefficients of $t^{2n}$:
\[
\frac{H_{2n}(0)}{(2n)!} = \frac{(-1)^n}{n!},
\]
so:
\[
\boxed{H_{2n}(0) = (-1)^n\frac{(2n)!}{n!}.}
\]

\subsection*{Part (c): $H_n'(0) = 2nH_{n-1}(0)$}

The Hermite polynomials satisfy the recurrence relation:
\[
H_n'(x) = 2nH_{n-1}(x).
\]

This can be derived from the generating function. Differentiating ($*$) with respect to $x$:
\[
\frac{\partial}{\partial x}e^{2xt-t^2} = 2t\,e^{2xt-t^2} = \sum_{n=0}^{\infty}H_n'(x)\frac{t^n}{n!}.
\]

But also:
\[
2t\,e^{2xt-t^2} = 2t\sum_{n=0}^{\infty}H_n(x)\frac{t^n}{n!} = \sum_{n=0}^{\infty}2H_n(x)\frac{t^{n+1}}{n!} = \sum_{n=1}^{\infty}\frac{2H_{n-1}(x)}{(n-1)!}\cdot t^n.
\]

Comparing coefficients of $t^n$:
\[
\frac{H_n'(x)}{n!} = \frac{2H_{n-1}(x)}{(n-1)!},
\]
giving:
\[
H_n'(x) = 2n\,H_{n-1}(x).
\]

Setting $x = 0$:
\[
\boxed{H_n'(0) = 2n\,H_{n-1}(0).}
\]

As a check: $H_2'(0) = 2\cdot 2\cdot H_1(0) = 4\cdot 0 = 0$. But $H_2(x) = 4x^2-2$, so $H_2'(x) = 8x$ and $H_2'(0) = 0$. Wait---let's check: $H_1(0) = 0$ (odd polynomial), and $H_2'(0) = 2\cdot 2 \cdot H_1(0) = 4\cdot 2\cdot 0 = 0$... actually $H_1(x) = 2x$, so $H_1(0) = 0$ and $H_2'(0) = 4\cdot 0 = 0$, but directly $H_2'(0) = 8\cdot 0 = 0$. Consistent.

Another check: $H_3'(0) = 6 H_2(0) = 6(-2) = -12$. Indeed $H_3(x) = 8x^3 - 12x$, so $H_3'(0) = -12$. $\checkmark$

\end{document}
